% TOP_PROBLEM: {"problemId": "problem.np-vs-ppoly", "canonicalTitle": "NP versus P/poly", "releaseRank": 12, "primaryDomain": "theoretical_computer_science", "primaryDomainLabel": "Theoretical computer science", "displayStatus": "open", "releaseStatus": "open"}
% !TeX program = lualatex
% Definition and sources only; this file is not an LLM solution attempt.
\documentclass[11pt]{article}
\usepackage[margin=25mm]{geometry}
\usepackage{fontspec}
\setmainfont{Latin Modern Roman}
\usepackage{amsmath,amssymb,mathtools}
\usepackage{unicode-math}
\setmathfont{Latin Modern Math}
\usepackage{xurl,hyperref}
\hypersetup{colorlinks=true,urlcolor=blue}
\setcounter{secnumdepth}{0}
\setlength{\parindent}{0pt}
\setlength{\parskip}{5pt}
\setlength{\emergencystretch}{3em}
\begin{document}
\section*{\texorpdfstring{NP versus P/poly}{NP versus P/poly}}\label{12}
\subsection{Definitions and mathematical statement}
A Boolean circuit is a finite directed acyclic computation with binary AND/OR gates, unary NOT gates, input bits, constants, and one output. Define $\mathsf{P/poly}$ by
\[
 L\in\mathsf{P/poly}\ \Longleftrightarrow\ \exists C,k\ \forall n\ge0\ \exists C_n:
 |C_n|\le C(n+1)^k,\quad C_n(x)=1_L(x)\ (|x|=n).
\]
All nodes are counted. No algorithm for generating $C_n$ is required. With $\mathsf{NP}$ defined by polynomially checkable witnesses, determine whether $\mathsf{NP}\subseteq\mathsf{P/poly}$. Different languages may have different size bounds.
\par\smallskip\textit{Formulation and concept references:} \href{https://theory.cs.princeton.edu/complexity/book.pdf}{[S1]}.

\subsection{Short English statement}
Can every efficiently verifiable decision problem be solved by small circuits when a different circuit may be chosen separately for each input length?
\subsection{Sources}
\begin{itemize}
\item[S1]Arora and Barak, Computational Complexity: A Modern Approach, web draft January8,2007, Definitions6.1–6.3 and Theorem6.13; Aaronson, P ?= NP, Conjecture25.
\url{https://theory.cs.princeton.edu/complexity/book.pdf}.
\textit{Formulation source.}
\item[S2]Boolean Circuits and P/poly, Cornell CS 6810, February 26, 2026.
\url{https://www.cs.cornell.edu/courses/cs6810/2026sp/lec10.pdf}.
\textit{Further reference; not a proof endorsement.}
\end{itemize}
\end{document}
